\section{Experimental Design}\label{sec:setup}
The experiment is a cross product of four axes --- encoder, input geometry, readout head and
dataset --- evaluated under exact protocols. \Cref{tab:encoders} lists the encoders,
\cref{tab:datasets} the datasets and protocols, and \cref{tab:coverage} which cells of the
product are measured, so that an absence is visible as an absence rather than as a blank.

\begin{table*}[t]
\centering
\caption{The encoders. Every checkpoint is frozen: no gradient reaches a backbone anywhere in
this study. \emph{fit} is how an image is made to match the input size --- \texttt{squash}
resizes to the target, \texttt{crop} short-side-resizes and centre-crops. C-RADIOv4 distils
SigLIP2-g, DINOv3 and a segmentation teacher, so rows 1--6 are students of rows 7--10's
families.}
\label{tab:encoders}
\small
\setlength{\tabcolsep}{4pt}
% generated by tools/gen_tables.py from results/runs — edit the runs, not this
\begin{tabular}{lllllp{0.26\linewidth}}
\toprule
encoder & pre-training & params (M) & input & fit & checkpoint \\
\midrule
C-RADIOv4-H & agglomerative & 653 & 224$\times$224 & squash & \texttt{\small torchhub:\allowbreak NVlabs/\allowbreak RADIO/\allowbreak c-radio\_\allowbreak v4-h} \\
C-RADIOv4-H & agglomerative & 653 & 256$\times$128 & squash & \texttt{\small torchhub:\allowbreak NVlabs/\allowbreak RADIO/\allowbreak c-radio\_\allowbreak v4-h} \\
C-RADIOv4-H & agglomerative & 653 & \textit{native} & squash & \texttt{\small torchhub:\allowbreak NVlabs/\allowbreak RADIO/\allowbreak c-radio\_\allowbreak v4-h} \\
C-RADIOv4-SO400M & agglomerative & 431 & 224$\times$224 & squash & \texttt{\small torchhub:\allowbreak NVlabs/\allowbreak RADIO/\allowbreak c-radio\_\allowbreak v4-so400m} \\
C-RADIOv4-SO400M & agglomerative & 431 & 256$\times$128 & squash & \texttt{\small torchhub:\allowbreak NVlabs/\allowbreak RADIO/\allowbreak c-radio\_\allowbreak v4-so400m} \\
C-RADIOv4-SO400M & agglomerative & 431 & \textit{native} & squash & \texttt{\small torchhub:\allowbreak NVlabs/\allowbreak RADIO/\allowbreak c-radio\_\allowbreak v4-so400m} \\
SigLIP2-g & language-aligned & 1163 & 256$\times$256 & squash & \texttt{\small timm:\allowbreak vit\_\allowbreak giantopt\_\allowbreak patch16\_\allowbreak siglip\_\allowbreak 256.v2\_\allowbreak webli} \\
SigLIP2-SO400M & language-aligned & 428 & 224$\times$224 & squash & \texttt{\small timm:\allowbreak vit\_\allowbreak so400m\_\allowbreak patch14\_\allowbreak siglip\_\allowbreak 224.v2\_\allowbreak webli} \\
DINOv3-H+ & self-supervised & 840 & 224$\times$224 & squash & \texttt{\small timm:\allowbreak vit\_\allowbreak huge\_\allowbreak plus\_\allowbreak patch16\_\allowbreak dinov3.lvd1689m} \\
DINOv3-L & self-supervised & 303 & 224$\times$224 & squash & \texttt{\small timm:\allowbreak vit\_\allowbreak large\_\allowbreak patch16\_\allowbreak dinov3.lvd1689m} \\
CLIP-B/16 & language-aligned & 87 & 224$\times$224 & squash & \texttt{\small timm:\allowbreak vit\_\allowbreak base\_\allowbreak patch16\_\allowbreak clip\_\allowbreak 224.openai} \\
CLIP-B/16 & language-aligned & 87 & 224$\times$224 & crop & \texttt{\small timm:\allowbreak vit\_\allowbreak base\_\allowbreak patch16\_\allowbreak clip\_\allowbreak 224.openai} \\
\bottomrule
\end{tabular}

\end{table*}

\subsection{Encoders}
\nEncoders{} backbones at \nEncoderConfigs{} encoder-resolution configurations span three
pre-training paradigms. \textbf{C-RADIOv4} (H, 653M; SO400M, 431M) is
agglomerative: one student distilled from SigLIP2-g, DINOv3-7B and a segmentation teacher
\cite{ranzinger2024amradio,heinrich2025radiov25}. \textbf{SigLIP2}
\cite{tschannen2025siglip2} appears at 1163M --- the literal teacher --- and at
428M, capacity-matched to C-RADIOv4-SO400M. \textbf{DINOv3}
\cite{simeoni2025dinov3} appears at 840M and 303M, bracketing C-RADIOv4-H's
653M from above and below; note that the true teacher is DINOv3-7B, which does not fit
on the card this study ran on, so the DINOv3 rows stand in for a teacher family rather than
being the teacher. \textbf{CLIP ViT-B/16} \cite{radford2021clip} is the small, well-understood
control, and is the one encoder present under both resize geometries.

Features are the encoder's pooled summary representation at \texttt{fp32}, cached to a
content-addressed store keyed on the encoder description and the manifest digest. C-RADIOv4 is
additionally run at $256\times128$ and at each image's native size, snapped to the model's
patch grid and capped at 512 pixels, which is the resolution axis of \cref{sec:geometry}.

\subsection{Readout heads}
A \emph{head} is a small map fitted on the frozen encoder's cached features and nothing else.
\nHeadRecipes{} recipes are reported, and \texttt{none} is one of them:
\begin{itemize}
  \item \textbf{none} --- the frozen encoder, ranked by cosine similarity. The zero-shot
        convention, and the row every other row is read against.
  \item \textbf{pca} --- a 512-dimensional PCA projection fitted on
        \texttt{market1501/train}. It sees the same \nHeadImages{} images as the supervised
        heads and none of the \nHeadIdentities{} identity labels. This is the control that
        separates \emph{supervision} from \emph{dimensionality}.
  \item \textbf{linear} --- a 512-dimensional affine map trained with softmax
        cross-entropy over the \nHeadIdentities{} training identities; the classifier is
        discarded and retrieval is scored on the projected features.
  \item \textbf{arcface} --- the same map trained with an additive angular margin
        \cite{deng2019arcface} ($m=0.3$, $s=30$).
\end{itemize}
All three fitted heads share one budget --- 100 epochs, Adam, lr $10^{-3}$, weight decay
$5\times10^{-4}$, batch 256, $\ell_2$-normalised inputs, seed 0 --- fixed on training accuracy
before any retrieval number was read. Fitting takes well under a minute per encoder on the
study's single GPU. Except where \cref{tab:fitsplit} says otherwise, every head is fitted on
\texttt{market1501/train} and applied unchanged everywhere, which makes exactly one column of
every results table in-domain supervised and every other column cross-domain transfer with no
target label ever seen. That asymmetry is the most important single fact about these tables, is
why the regime is named in \cref{tab:datasets}, and is measured rather than assumed in
\cref{sec:heads}, where the same recipes are refitted on a second training split.

\begin{table*}[t]
\centering
\caption{Datasets and protocols. \emph{regime} is relative to the heads, all of which are
fitted on \texttt{market1501/train}; only Market-1501 is in-domain. Query and gallery sizes are
read from the run records, so a cross-dataset comparison of absolute scores can be checked
against the search size that produced it.}
\label{tab:datasets}
\scriptsize
\setlength{\tabcolsep}{3pt}
% generated by tools/gen_tables.py from results/runs — edit the runs, not this
\begin{tabular}{llllrr}
\toprule
dataset & object & regime & protocol & queries & gallery \\
\midrule
Market-1501 & person & in-domain & \texttt{market1501/official@1} & 3,368 & 19,732 \\
MSMT17 & person & transfer & \texttt{msmt17/official@1} & 11,659 & 82,161 \\
CUHK03-NP & person & transfer & \texttt{cuhk03/detected-767@1} & 1,400 & 5,332 \\
Occluded-REID & person & transfer & \texttt{occluded-reid/occluded-vs-whole@1} & 1,000 & 1,000 \\
CCVID & person & transfer & \texttt{ccvid/tracklet-cloth-changing@1} & 834 & 1,074 \\
CCVID & person & transfer & \texttt{ccvid/tracklet@1} & 834 & 1,074 \\
VRIC & vehicle & transfer & \texttt{vric/official@1} & 2,811 & 2,811 \\
VRAI & vehicle & transfer & \texttt{vrai/train-cross-camera@1} & 6,302 & 32,338 \\
\bottomrule
\end{tabular}

\end{table*}

\subsection{Datasets and protocols}
\nDatasets{} datasets under \nProtocols{} protocols cover five regimes of difficulty.
\textbf{Market-1501} \cite{zheng2015scalable} is the source domain and the only in-domain
column. \textbf{MSMT17} \cite{wei2018msmt} is the large-scale person transfer target.
\textbf{CUHK03-NP} \cite{zhong2017reranking}, detected boxes under the 767/700 split, is the
degraded-detector regime and the one closest to what a real detector hands a ReID model.
\textbf{Occluded-REID} \cite{zhuo2018occluded} ships no standard split and labels no cameras,
so the whole set is used, occluded probes against whole-body gallery, and the same-camera
exclusion every Market number depends on does not exist there. \textbf{CCVID}
\cite{gu2022ccvid} is video and cloth-changing; frames are thinned to eight per tracklet before
encoding and mean-pooled into one vector per tracklet afterwards, and both protocols --- with
and without the cloth-changing restriction --- are reported. \textbf{VRIC} \cite{kanaci2018vric}
and \textbf{VRAI} \cite{wang2019vrai} are vehicles, the latter aerial. VRAI's test identities
are withheld by its release, so the only protocol that scores offline runs over its training
split, querying the first frame of each camera-1 trajectory against every camera-2 training
image; that is a legitimate zero-shot number for an encoder that never saw VRAI, and it is
sound for the head rows here \emph{only} because no head in this study has ever seen a VRAI
image.

We do not use DukeMTMC or any dataset derived from it. The exclusion is enforced in code with
no override.

\begin{table*}[t]
\centering
\caption{Coverage: run records per (encoder, dataset) pair, four per cell being the four heads
and eight where the dataset ships two protocols. The MSMT17 column is the study's one
systematic gap and is discussed in \cref{sec:limitations}.}
\label{tab:coverage}
\scriptsize
\setlength{\tabcolsep}{3pt}
% generated by tools/gen_tables.py from results/runs — edit the runs, not this
\begin{tabular}{llrrrrrrr}
\toprule
encoder & input & Market & MSMT17 & CUHK03 & Occ-REID & CCVID & VRIC & VRAI \\
\midrule
C-RADIOv4-H & 224$\times$224 & 7 & 7 & 7 & 7 & 7 & 7 & 7 \\
C-RADIOv4-H & 256$\times$128 & 7 & 7 & 7 & 7 & 7 & 7 & 7 \\
C-RADIOv4-H & \textit{native} & 7 & 7 & 7 & 7 & 7 & 7 & 7 \\
C-RADIOv4-SO400M & 224$\times$224 & 7 & 7 & 7 & 7 & 7 & 7 & 7 \\
C-RADIOv4-SO400M & 256$\times$128 & 7 & 7 & 7 & 7 & 7 & 7 & 7 \\
C-RADIOv4-SO400M & \textit{native} & 7 & 7 & 7 & 7 & 7 & 7 & 7 \\
SigLIP2-g & 256$\times$256 & 7 & 7 & 7 & 7 & 7 & 7 & 7 \\
SigLIP2-SO400M & 224$\times$224 & 7 & 7 & 7 & 7 & 7 & 7 & 7 \\
DINOv3-H+ & 224$\times$224 & 7 & 7 & 7 & 7 & 7 & 7 & 7 \\
DINOv3-L & 224$\times$224 & 7 & 7 & 7 & 7 & 7 & 7 & 7 \\
CLIP-B/16 (squash) & 224$\times$224 & 7 & 7 & 7 & 7 & 7 & 7 & 7 \\
CLIP-B/16 (crop) & 224$\times$224 & 7 & 7 & 7 & 7 & 7 & 7 & 7 \\
\bottomrule
\end{tabular}

\end{table*}

\subsection{Metrics}
We report mAP and Rank-1 throughout, and mINP \cite{ye2022deep} where the claim is about the
hardest true match rather than the average one. Cross-dataset comparison of absolute values is
not meaningful here and we do not make it: Occluded-REID searches roughly a thousand images,
Market-1501 searches nearly twenty thousand, and VRAI over thirty thousand, so a gallery an
order of magnitude larger accounts for much of the difference between those columns before any
question of difficulty arises. Every claim in this paper is a within-column comparison, where
all rows share a protocol digest and a manifest digest.

\subsection{Noise}
Refitting the lead encoder's ArcFace head at three seeds and rescoring gives a standard
deviation of $0.0002$ mAP on Market-1501, $0.0050$ on Occluded-REID --- the noisiest set here,
with its thousand queries --- and $0.0009$ on VRAI. We treat differences below roughly twice
those values as ties and say so where it matters; several comparisons in \cref{sec:ablation}
turn on this.
